% from ega4-18-fr.tex
\oldpage[IV-4]{215}

Nous procéderons par \emph{récurrence sur} $m$, supposant donc (19.7.1.6) vérifiée lorsqu'on remplace $m$ par un entier $m'<m$. Nous considérerons d'autre part, pour $m$ fixé et pour un entier $d\geq0$, la relation
\begin{equation}
  R\cap(\mathfrak K^mN+\mathfrak L^{m+1}N)\subset(R\cap\mathfrak K^dN\cap\mathfrak L^mN)+\mathfrak L^{m+1}N.
  \tag{$*_d$}
\end{equation}

Il est clair qu'il suffit de démontrer $(*_d)$ pour $d=m$, et d'autre part $(*_0)$ est trivialement vraie. Nous prouverons $(*_d)$ par \emph{récurrence sur} $d$; autrement dit, nous supposons, pour un $d\geq0$ fixé (avec $d<m$), que $(*_d)$ est vraie, et nous voulons prouver
\begin{equation}
  R\cap(\mathfrak K^mN+\mathfrak L^{m+1}N)\subset(R\cap\mathfrak K^{d+1}N\cap\mathfrak L^mN)+\mathfrak L^{m+1}N.
  \tag{$*_{d+1}$}
\end{equation}

Considérons pour cela, pour un $h\geq0$, la relation
\begin{equation}
  R\cap(\mathfrak K^mN+\mathfrak L^{m+1}N)\subset(R\cap(\mathfrak K^{d+1}N+\mathfrak K^d\mathfrak L^hN)\cap\mathfrak L^mN)+\mathfrak L^{m+1}N
  \tag{$**_{d+1,h}$}
\end{equation}

L'hypothèse $(*_d)$ implique que $(**_{d+1,0})$ est vraie. Nous prouverons par \emph{récurrence sur} $h$ que $(**_{d+1,h})$ est vraie pour \emph{tout} $h\geq0$. Mais on a le

\begin{lemma}[19.7.1.7]
\label{IV.19.7.1.7}
Soient $A$ un anneau noethérien, $N$ un $A$-module de type fini, $E$, $F$ deux sous-modules de $N$, $\mathfrak L$ un idéal de $A$. Pour tout $k>0$ il existe $h>0$ tel que $E\cap(F+\mathfrak L^hN)\subset(E\cap F)+\mathfrak L^kN$.
\end{lemma}

En effet, soit $\varphi:N\to N/(E\cap F)=N_1$ l'homomorphisme canonique, et posons $E_1=\varphi(E)$, $F_1=\varphi(F)$, de sorte que $E_1\cap F_1=0$ et $\varphi(E\cap(F+\mathfrak L^hN))=E_1\cap(F_1+\mathfrak L^hN_1)$. On sait ($\mathbf{0}_{\mathrm I}$, 7.3.2) qu'il existe $h$ tel que $(E_1+F_1)\cap\mathfrak L^hN_1\subset\mathfrak L^k(E_1+F_1)$; si $h$ est ainsi choisi et si $x_1\in E_1$ est tel qu'il existe $y_1\in F_1$ pour lequel $x_1\setmin y_1\in\mathfrak L^hN_1$, on en déduit d'après ce qui précède qu'il existe $u_1\in\mathfrak L^kE_1$ et $v_1\in\mathfrak L^kF_1$ tels que $x_1=u_1$ et $y_1=v_1$; donc $E_1\cap(F_1+\mathfrak L^hN_1)\subset\mathfrak L^kN_1$, ce qui prouve le lemme.

Il suffit alors d'appliquer ce lemme avec $E=R\cap\mathfrak L^mN$, $F=\mathfrak K^{d+1}N$ et de prendre $k=m+1$ pour en déduire qu'avec la valeur de $h$ correspondante, $(**_{d+1,h})$ entraîne $(*_{d+1})$.

Nous supposons donc dans tout ce qui suit que $(**_{d+1,h})$ \emph{est vérifiée}.

\begin{env}[19.7.1.8]
\label{IV.19.7.1.8}
\emph{Démonstration de} $(**_{d+1,h+1})$: \emph{Premier cas}: $h\leq m\setmin d$.
\end{env}

Partons donc d'un élément
\begin{env}[19.7.1.9]
\label{IV.19.7.1.9}
\[
  g\in R\cap(\mathfrak K^{d+1}N+\mathfrak K^d\mathfrak L^hN)\cap\mathfrak L^mN
\]
\end{env}
congru mod. $\mathfrak L^{m+1}N$ à un élément $g_1\in R\cap(\mathfrak K^mN+\mathfrak L^{m+1}N)$; on a donc
\begin{env}[19.7.1.10]
\label{IV.19.7.1.10}
\[
  g\in\mathfrak K^mN+\mathfrak L^{m+1}N.
\]
\end{env}

Il suffira de prouver que $g$ est aussi congru mod. $\mathfrak L^{m+1}N$ à un élément
\[
  g_2\in R\cap(\mathfrak K^{d+1}N+\mathfrak K^d\mathfrak L^{h+1}N),
\]
car cela entraînera $g_2\in\mathfrak L^mN$, puisque l'on aura $g\in\mathfrak L^mN$ et $g_2\setmin g\in\mathfrak L^{m+1}N$. Dans le cas $h\leq m\setmin d$ que nous considérons, il suffit de démontrer la relation
\begin{env}[19.7.1.11]
\label{IV.19.7.1.11}
\[
  (\mathfrak K^mN+\mathfrak L^{m+1}N)\cap\mathfrak K^dN\subset\mathfrak L^{m-d+1}\mathfrak K^dN+\mathfrak K^{d+1}N
\]
\end{env}
